\section{Model Selection}
In Model Selection, we seek to choose the best hyperparameter / learning
algorithm / model for a particular problem. We are primarily concerned with
ranking a set of models, rather than estimating the performance of a particular
model.
\subsubsection{Measures for Performance}
The measure of performance we are ultimately interested in is \textbf{Generalisation Loss}:
\[\mathbb{E}_\mathcal{D}\left[ \mathcal{E}\left( f(\mathbf{\mathcal{X}}) ,\mathcal{Y}\right) \right].\] 
This cannot be directly measured, and empirical test set loss,
$\mathbb{E}_\mathcal{S_{Test}}\left[ \mathcal{E}\left( f(\mathbf{\mathcal{X}})
,\mathcal{Y}\right) \right]$ can act as a reasonable estimator (but not for
model selection). Empirical Training Set Loss cannot be used as well, as it
will lead to overfitting.\\
Hence, Empirical Risk Minimisation alone is not adequate to perform model
selection.

\subsection{Validation Techniques}
We seek to estimate the generalisation loss through a \textbf{Validation Set
Loss}. We split the training data into three sets:
\begin{itemize}
  \item Training: To train each possible model
  \item Validation: To select the `best' model
  \item Test: To assess the performance of that model
\end{itemize}
\begin{figure}[h!]
  \centering
  \includegraphics[width=0.8\textwidth]{validation}
\end{figure}
\subsubsection{Validation Loss Bound}
Applying the Hoeffding's Inequality to i.i.d. $\mathcal{Z}_i=\mathbb{I}\left[ \mathcal{Y}^{(i)}\neq f\left( \mathcal{\mathbf{X}}^{(i)} \right) \right]$:
\[P\left( E_\mathcal{D}\left[ \mathbb{I}\left[ \mathcal{Y}\neq f\left( \mathcal{\mathbf{X}} \right)  \right] \right] \ge \frac{1}{n_v}\sum_{i=1}^{n_v} \mathbb{I}\left[ y^{(i)}\neq f(\mathbf{x}^{(i)}) \right] + \epsilon\right)\le e^{-2n_v\epsilon^{2}}.\] 
\[P\left( E_\mathcal{D}\left[ \mathbb{I}\left[ \mathcal{Y}\neq f\left( \mathcal{\mathbf{X}} \right)  \right] \right] \ge E_{S_v} \left[ \mathbb{I}\left[ y^{(i)}\neq f(\mathbf{x}^{(i)}) \right] \right] + \epsilon\right)\le e^{-2n_v\epsilon^{2}}.\] 
Let $\delta=e^{-2n_v\epsilon^{2}}\implies\epsilon=\sqrt{\frac{\ln\left( \frac{1}{\delta} \right) }{2n_v}}$.
\[P\left( E_\mathcal{D}\left[ \mathbb{I}\left[ \mathcal{Y}\neq f\left( \mathcal{\mathbf{X}} \right)  \right] \right] < E_{S_v} \left[ \mathbb{I}\left[ \mathcal{Y}\neq f(\mathcal{\mathbf{X}}) \right] \right] + \sqrt{\frac{\ln\left( \frac{1}{\delta} \right) }{2n_v}}\right)> (1-\delta).\] 
That is, with probability greater than $(1-\delta)$, we state that:
\[L(\mathcal{E}, \mathcal{D}, f)\le L(\mathcal{E}, \mathcal{S}_v, f)+\sqrt{\frac{\ln\left( \frac{1}{\delta} \right) }{2n_v}}.\] 
Using the Union Bound, we can show that a bound holds for all members in a finite hypothesis class
$\mathcal{F}$:
Let $A,B$ be the events:
\[A=\left\{ L(\mathcal{E}, \mathcal{D}, f^{A})\le L(\mathcal{E}, \mathcal{S}_v, f^{A})+\sqrt{\frac{\ln\left( \frac{1}{\delta'} \right) }{2n_v}} \right\}.\] 
\[B=\left\{ L(\mathcal{E}, \mathcal{D}, f^{B})\le L(\mathcal{E}, \mathcal{S}_v, f^{B})+\sqrt{\frac{\ln\left( \frac{1}{\delta'} \right) }{2n_v}} \right\}.\] 
We have $P(A)=P(B)<\delta'$.
\begin{remark}
  \[P(\neg A\wedge\neg B)\ge 1-(P(A) + P(B)).\]
\end{remark}
Combining these, $\forall f \in \left\{ f^{A},f^{B} \right\}$:
\[P\left( L(\mathcal{E}, \mathcal{D}, f)\le L(\mathcal{E}, \mathcal{S}_v, f)+\sqrt{\frac{\ln\left( \frac{1}{\delta'} \right) }{2n_v}} \right)\ge 1-2\delta'.\] 
Let $\delta=2\delta'$, and we have:
\[P\left( L(\mathcal{E}, \mathcal{D}, f)\le L(\mathcal{E}, \mathcal{S}_v, f)+\sqrt{\frac{\ln\left( \frac{2}{\delta} \right) }{2n_v}} \right)\ge 1-\delta.\] 
In general, $\forall f\in\mathcal{F}$:
\[P\left( L(\mathcal{E}, \mathcal{D}, f)\le L(\mathcal{E}, \mathcal{S}_v, f)+\sqrt{\frac{\ln\left( \frac{\left| {\mathcal{F}} \right|}{\delta} \right) }{2n_v}} \right)\ge 1-\delta.\] 
The validation is a good estimate provided $\mathcal{F}$ is not too large, and
$n_v$ is not too small. Otherwise, the validation and generalisation losses
will vary further as we begin to overfit.

\subsection{Cross Validation}
To retain most of our data for training, we perform validation while still
maintaining the training data set, by spliting the training data into $k$
folds. The learning algorithm is run $k$ times for each model, each time using
all the folds but one as the training set, and the remaining fold as the
validation set.
\begin{figure}[h!]
  \centering
  \includegraphics[width=0.8\textwidth]{cross-validation}
\end{figure}
The cross validation loss is the performance averaged across all $k$ folds:
\[L_{{CV}_k}\left( \mathcal{E}, \mathcal{S}, f \right) =\frac{1}{k}\sum_{i=1}^{k} L_{\mathcal{S}\backslash\mathcal{S}_k}\left( \mathcal{E},\mathcal{S}\backslash\mathcal{S}_i ,f \right).\] 
\subsubsection{$k$ selection}
$k=n$ is known as the Leave One Out cross validation, maximising the training
set for each fold, is computation expensive, and fails to give a good estimate
of the generalisation loss.\\
Empirical analysis suggests $k=5$ or $10$.
\subsubsection{Limitations}
There is not yet a rigorous understanding of why cross validation works well,
and in certain circumstances should be treated with care.
\begin{example}[Time Series Data]
  Requires sequential cross validation, due to similar neighbouring points.
\end{example}

\subsection{PAC Approach}
In the PAC (Probably Approximately Correct) approach, we formulate a
probabilistic `worst-case' bound using the empirical training loss and some
complexity penalty (taking into account $\left| \mathcal{F} \right| $). We then
perform based on the model that provides the tightest bound.\\\\
We obtain these bounds using the Uniform Convergence approach, by first seeking
a finite sample concentration inequality independent of $\mathcal{D}$ to one
function $f$, then applying it uniformly across all $f \in \mathcal{F}$ via
Union Bound or other measure of functional complexity.\\\\
We consider a probability of success $(1-\delta)$, an approximation loss target
$\epsilon$, and a data generating process $\mathcal{D}$.

\subsubsection{Consistent Learner}
For a consistent learning algorithm that only outputs hypotheses consistent
with the training data ($E_\mathcal{S}\left[
\mathcal{E}(f(\mathcal{X}),\mathcal{Y}) \right]=0$), for
$0<\epsilon<1,0<\delta<1$ and a particular $f \in \mathcal{F}\wedge
E_\mathcal{D}\left[ \mathcal{E}\left( f(\mathcal{\mathbf{X}}),\mathcal{Y}
\right) \right]>\epsilon$:
\[P(f\text{ is consistent with 1 training ex.})<(1-\epsilon)<e^{-\epsilon}.\] 
\[P(f\text{ is consistent with $n$ training ex.})<{(1-\epsilon)}^{n}<e^{-\epsilon n}.\] 
Applying the Union Bound $\forall f \in \mathcal{F}$, we obtain:
\[P\left( E_\mathcal{D}\left[ \mathcal{E}\left( f(\mathbf{X}),\mathcal{Y} \right) \right]>\epsilon, \text{ for at least one $f \in \mathcal{F}$}\right)<\left| \mathcal{F} \right|e^{-\epsilon n}.\] 
\begin{theorem}[PAC Bound: Finite $\mathcal{F}$, Consistent Learner]
  For all $0<\epsilon<1$, $0<\delta<1$, all data generating distributions $\mathcal{D}$, with
  probability of at least $(1-\delta)$:
  \[\forall f \in \mathcal{F}\quad E_\mathcal{D}\left[ \mathcal{E}\left( f(\mathbf{\mathcal{X}}),\mathcal{Y} \right) \right]\le \frac{1}{n}\ln\left( \frac{\left| \mathcal{F} \right|}{\delta}  \right).\] 
\end{theorem}

\subsubsection{Agnostic Learner}
For a learning algorithm which is not restricted in the hypotheses:
Applying the Hoeffding's Inequality to i.i.d. $\mathcal{Z}_i=\mathbb{I}\left[ \mathcal{Y}^{(i)}\neq f\left( \mathcal{\mathbf{X}}^{(i)} \right) \right]$:
\[P\left( E_\mathcal{D}\left[ \mathbb{I}\left[ \mathcal{Y}\neq f\left( \mathcal{\mathbf{X}} \right)  \right] \right] \ge \frac{1}{n}\sum_{i=1}^{n} \mathbb{I}\left[ y^{(i)}\neq f(\mathbf{x}^{(i)}) \right] + \epsilon\right)\le e^{-2n\epsilon^{2}}.\] 
\[P\left( E_\mathcal{D}\left[ \mathbb{I}\left[ \mathcal{Y}\neq f\left( \mathcal{\mathbf{X}} \right)  \right] \right] \ge E_{S} \left[ \mathbb{I}\left[ \mathcal{Y}\neq f(\mathcal{\mathbf{X}}) \right] \right] + \epsilon \right)> e^{-2n\epsilon^{2}}.\] 
Similarly, we applying the Union Bound to obtain:
\[P\left( E_\mathcal{D}\left[ \mathbb{I}\left[ \mathcal{Y}\neq f\left( \mathcal{\mathbf{X}} \right)  \right] \right] \ge E_{S} \left[ \mathbb{I}\left[ \mathcal{Y}\neq f(\mathcal{\mathbf{X}}) \right] \right] + \epsilon\text{, for at least one }f \in \mathcal{F}\right)\le \left| \mathcal{F} \right| e^{-2n\epsilon^{2}}.\] 
\begin{theorem}[PAC Bound: Finite $\mathcal{F}$, Agnostic Learner]
  For all $0<\epsilon<1$, $0<\delta<1$, all data generating distributions $\mathcal{D}$, with
  probability of at least $(1-\delta)$:
\[\forall f \in \mathcal{F}\quad E_\mathcal{D}\left[ \mathbb{I}\left[ \mathcal{Y}\neq f\left( \mathcal{\mathbf{X}} \right)  \right] \right] \le E_{S} \left[ \mathbb{I}\left[ \mathcal{Y}\neq f(\mathcal{\mathbf{X}}) \right] \right] + \sqrt{\frac{1}{2n}\ln\left( \frac{\left| \mathcal{F} \right|}{\delta}  \right) }\] 
\end{theorem}

\subsubsection{Infinite $\mathcal{F}$}
\begin{theorem}[PAC Bound: Infinite $\mathcal{F}$, Agnostic Learner]
  For all $0<\epsilon<1$, $0<\delta<1$, all data generating distributions $\mathcal{D}$, with
  probability of at least $(1-\delta)$:
\[\forall f \in \mathcal{F}\quad E_\mathcal{D}\left[ \mathbb{I}\left[ \mathcal{Y}\neq f\left( \mathcal{\mathbf{X}} \right)  \right] \right] \le E_{S} \left[ \mathbb{I}\left[ \mathcal{Y}\neq f(\mathcal{\mathbf{X}}) \right] \right] + 2\mathbb{R}\left( \mathcal{F}\circ\mathcal{S} \right) + 4\sqrt{\frac{2}{n}\ln\left( \frac{4}{\delta}  \right) }\] 
where $\mathbb{R}\left( \mathcal{F}\circ\mathcal{S} \right)$ denotes the Rademacher complexity of $\mathcal{F}\circ\mathcal{S}={\left\{ f(\mathbf{x}^{(i)})|f \in \mathcal{F} \right\}}_{i=1}^{n}$
\end{theorem}

\subsubsection{Advantages and Limitations}
The PAC approach is not limited to misclassification loss, but to a variety of
other loss functions. It allows comparison with training data alone while
avoiding overfitting.\\

However, the bounds are often loose, and are poor estimators of generalisation
loss, and tend to underperform cross validation.
